% !TEX TS-program = lualatex+se
% !TEX encoding = UTF-8 Unicode

%\errorcontextlines=0
\documentclass[17pt, a4paper]{extbook}
% -----------------------------------------------------------------------------------
% ------------------------- USTAWIENIA STRONY--------------------------
% -----------------------------------------------------------------------------------
\usepackage{geometry}
\geometry{
	left=15mm,
	right=15mm,
	top=15mm,
	bottom=15mm,
	marginparwidth=0mm,
	marginparsep=0mm
	}
%\newcommand{\wstawplik}[1]{\input{/home/w5rdshat2-h/Documents/Latex/include/#1}}%skraca kod
%\usepackage[letterpaper,showframe]{geometry}	
% usuwanie nagłówków/stopek z\,dokumentu
%\usepackage[indent=0pt]{parskip} % ustawianie paragrafów
%\setstretch{1.5} % półtorej linii odległości między wierszami
%\usepackage{titlesec} % ustawianie nagłówków
%%\usepackage{fullpage} % redukcja marginesów
%\usepackage{setspace} % odstępy między liniami
%\usepackage{enumitem}% listy numerowane
\usepackage{soul}% dodaje \ul
%\usepackage{ulem}% dodaje \uline

% -----------------------------------------------------------------------------------
% ------------------------------ GREGORIO -----------------------------------
% -----------------------------------------------------------------------------------
%\usepackage{gregoriotex}
%\gresetlinecolor{gregoriocolor}% Make staff lines red; remove for black:
%%\gresetgregoriofont{granapadano} %{gregorio}
%\gresetheadercapture{annotation}{greannotation}{string}% odczytuje adnotacje nad inicjałem jako ciąg, co pozwala na jego renderowanie
%\gresetheadercapture{commentary}{grecommentary}{string}% odczytuje komentarz jako ciąg, co pozwala na jego renderowanie
%\grechangedim{commentaryraise}{0.3cm}{scalable} % podnosi komentarz wystarczająco ponad górną linię
%\grechangestyle{annotation}{\color{black}\small}% format adnotacji
%\grechangestyle{commentary}{\color{gregoriocolor}\small}% format komentarza
%\grechangestyle{abovelinestext}{\color{gregoriocolor}\small}% format napisu nad liniami
%\grechangestyle{initial}{\fontsize{36}{36}\selectfont} % zmniejsza wielkość inicjału do bardziej proporcjonalnego rozmiaru (=36pt)
%\grechangedim{intersyllablespacenotes}{0.05cm}{scalable}% Minimalna odległość między nutami różnych sylab.
%% Set the space around the initial:
%% See http://gregorio-project.github.io/gregoriotex/details.html for more details and options
%\grechangedim{beforeinitialshift}{0.5mm}{scalable}
%\grechangedim{afterinitialshift}{1.0mm}{scalable}
%\grechangestaffsize{25}%rozmiar wyniku GregorioTEX

% -----------------------------------------------------------------------------------
% ---------------------- CZCIONKI DOKUMENTU ------------------------
% -----------------------------------------------------------------------------------
\usepackage{fontspec}
%\setmainfont{FreeSerif}
%\setsansfont{FreeSans}
%\setmonofont{FreeMono}
%\usepackage{ebgaramond}
\setmainfont{EB Garamond}
\defaultfontfeatures{Ligatures=TeX}
\usepackage{polyglossia}
%\usepackage{ecclesiastic}% This package works only with pdfLaTeX 
%\usepackage[autostyle=true]{csquotes}%generuje odpowiednie cudzysłowy dla danego języka

\setdefaultlanguage[variant=classic, hyphenation=liturgical,babelshorthands = true]{latin}
\setotherlanguage{polish}

%\usepackage[osf,p]{libertine}
%\setmainfont{Libertinus Serif}
%\setmainfont{TeX Gyre Termes}
%\setmainfont{TeX Gyre Heros}
%\setmainfont{Calluna Regular}

\begin{document}
{bez deklaracji:}\\
A-Á-Æ-Ǽ-B-C-D-E-É-F-G-H-I-Í-J-K-L-M-N-O-Ó-Œ-Œ́-P-Q-R-S-T-U-Ú-V-V́-W-X-Y-Ý-Z\\
a-á-æ-ǽ-b-c-d-e-é-f-g-h-i-í-j-k-l-m-n-o-ó-œ-œ́-p-q-r-s-t-u-ú-v-w-x-y-ý-z\\
'a 'e 'i 'o 'u 'y 'A 'E 'O 'U 'Y "ae "Ae "AE 'ae 'Ae 'AE "< ">\\
c\oe{}na - C\OE{}NA, tu\ae - TU\AE \\
\textit{A-Á-Æ-Ǽ-B-C-D-E-É-F-G-H-I-Í-J-K-L-M-N-O-Ó-Œ-Œ́-P-Q-R-S-T-U-Ú-V-V́-W-X-Y-Ý-Z\\
a-á-æ-ǽ-b-c-d-e-é-f-g-h-i-í-j-k-l-m-n-o-ó-œ-œ́-p-q-r-s-t-u-ú-v-w-x-y-ý-z\\
'a 'e 'i 'o 'u 'y 'A 'E 'O 'U 'Y "ae "Ae "AE 'ae 'Ae 'AE "< ">\\
c\oe{}na - C\OE{}NA, tu\ae - TU\AE \\}
% classic
{variant=classic, hyphenation=liturgical:}\\
\begin{latin}
A-Á-Æ-Ǽ-B-C-D-E-É-F-G-H-I-Í-J-K-L-M-N-O-Ó-Œ-Œ́-P-Q-R-S-T-U-Ú-V-V́-W-X-Y-Ý-Z\\
a-á-æ-ǽ-b-c-d-e-é-f-g-h-i-í-j-k-l-m-n-o-ó-œ-œ́-p-q-r-s-t-u-ú-v-w-x-y-ý-z\\
'a 'e 'i 'o 'u 'y 'A 'E 'O 'U 'Y "ae "Ae "AE 'ae 'Ae 'AE "< ">\\
c\oe{}na - C\OE{}NA, tu\ae - TU\AE \\
\end{latin}
% medieval
{variant=medieval, hyphenation=liturgical:}\\
\begin{latin}[variant=medieval, hyphenation=liturgical,babelshorthands = true]
A-Á-Æ-Ǽ-B-C-D-E-É-F-G-H-I-Í-J-K-L-M-N-O-Ó-Œ-Œ́-P-Q-R-S-T-U-Ú-V-V́-W-X-Y-Ý-Z\\
a-á-æ-ǽ-b-c-d-e-é-f-g-h-i-í-j-k-l-m-n-o-ó-œ-œ́-p-q-r-s-t-u-ú-v-w-x-y-ý-z\\
'a 'e 'i 'o 'u 'y 'A 'E 'O 'U 'Y "ae "Ae "AE 'ae 'Ae 'AE "< ">\\
c\oe{}na - C\OE{}NA, tu\ae - TU\AE \\
\end{latin}
% modern Æ (lowercase: æ
{[variant=modern, hyphenation=liturgical:}\\
\begin{latin}[variant=modern, hyphenation=liturgical,babelshorthands = true]
A-Á-Æ-Ǽ-B-C-D-E-É-F-G-H-I-Í-J-K-L-M-N-O-Ó-Œ-Œ́-P-Q-R-S-T-U-Ú-V-V́-W-X-Y-Ý-Z\\
a-á-æ-ǽ-b-c-d-e-é-f-g-h-i-í-j-k-l-m-n-o-ó-œ-œ́-p-q-r-s-t-u-ú-v-w-x-y-ý-z\\
'a 'e 'i 'o 'u 'y 'A 'E 'O 'U 'Y "ae "Ae "AE 'ae 'Ae 'AE "< ">\\
c\oe{}na - C\OE{}NA, tu\ae - TU\AE \\
\end{latin}
% ecclesiastic
{variant=ecclesiastic, hyphenation=liturgical:}\\
\begin{latin}[variant=ecclesiastic, hyphenation=liturgical,babelshorthands = true]
A-Á-Æ-Ǽ-B-C-D-E-É-F-G-H-I-Í-J-K-L-M-N-O-Ó-Œ-Œ́-P-Q-R-S-T-U-Ú-V-V́-W-X-Y-Ý-Z\\
a-á-æ-ǽ-b-c-d-e-é-f-g-h-i-í-j-k-l-m-n-o-ó-œ-œ́-p-q-r-s-t-u-ú-v-w-x-y-ý-z\\
'a 'e 'i 'o 'u 'y 'A 'E 'O 'U 'Y "ae "Ae "AE 'ae 'Ae 'AE "< ">\\
c\oe{}na - C\OE{}NA, tu\ae - TU\AE \\
\end{latin}
% polski
{polski:}\\
\begin{polish}
A-Ą-B-C-Ć-D-E-Ę-F-G-H-I-J-K-L-Ł-M-N-O-Ó-P-Q-R-S-Ś-T-U-V-W-X-Y-Z-Ź-Ż\\
a-ą-b-c-ć-d-e-ę-f-g-h-i-j-k-l-ł-m-n-o-ó-p-q-r-s-ś-t-u-v-w-x-y-z-ź-ż\\
\textbackslash{}textdagger - \textdagger\\
\textbackslash{}textdaggerdbl - \textdaggerdbl\\
\end{polish}
\end{document}